
\section{Tutorial}
\label{sec:manual-tutorial}

Welcome to CubedOS!

\subsection{Introduction}
\label{sec:manual-tutorial-introduction}

This document is for new CubedOS developers who want to learn how to build CubedOS applications.
Here we will walk you through the process of creating a simple application and then refer you to
the \filename{samples} folder in the CubedOS GitHub repository for additional examples.

CubedOS applications consist of one or more \newterm{domains} where each domain consists of one
or more \newterm{modules}. Simple (and even many not-so-simple) applications are built as a
single domain. Multi-domain applications are outside the scope of this tutorial, but the
Networking sample shows an example of such an application.

Each module in a domain communicates with other modules by passing messages. Sometimes modules
act as servers; they receive request messages, process those requests, and return reply messages
to the requester. Sometimes modules act as clients; they send request messages to servers,
receive the replies, and do other processing. It is normal for modules to play both roles at
different times; they receive request messages and send requests of their own as they process
their incoming requests.

CubedOS encourages a client/server application architecture, or, more generally, a distributed
application architecture with messages being sent between modules freely. However, the
Request/Reply terminology is conventionally used throughout. In some cases ``requests'' are sent
even when no reply is expected, and in other cases ``replies'' are sent unsolicited.

Normally modules are written in SPARK/Ada as a package hierarchy. The top-level, parent package
has the same name as the module. Child packages are used to implement the various subsystems
within a module. Some child packages have conventional names, as we will describe later. This
precise organization is not strictly required, but it is recommended and certain aspects of the
system assume this organization. It is possible, in principle, to write CubedOS modules in other
languages such as full Ada or even C. We do not discuss how to do that in this tutorial.

Communication in CubedOS is point-to-point, asynchronous, and unreliable. There is no built-in
mechanism for broadcasting or multicasting messages; doing so requires writing code that
explicitly sends multiple point-to-point messages as needed. There is, however, a
publish/subscribe core module that offers a ``channel'' abstraction. Modules can subscribe to
particular channels and receive messages that are published to those channels by other modules.
This decouples senders and receivers so they are not aware of each other and provides a form of
multicasting.

Messages are asynchronous in the sense that the sender is allowed to continue before the message
has been received. The send operation only queues the message for later delivery. After a
successful send, the sending module can reuse any data structures it populated for the outgoing
message. Every module has at least one task, called the \newterm{module main task} that receives
and processes messages. Modules may have other internal tasks for other purposes. Multiple
modules execute in parallel. It is possible, even common, for senders to send multiple requests,
possibly to different modules, before getting any replies, and replies may arrive in a different
order from which the requests were sent. It is up to you to coordinate this activity in your
modules, although CubedOS provides some assistance, as we will discuss.

Messages are unreliable in the sense that final delivery is not guaranteed. Senders can elect to
receive a status code to inform them if their messages are dropped by the message passing
infrastructure, but even this status code does not indicate if the message was actually received
by the destination module (for example, if the destination module is stuck or crashed or
unreachable over a network). If a server module wishes to indicate an error in a request, or any
other kind of failure, it must send an explicit reply message saying so.

Every CubedOS domain contains an array of ``mailboxes'' that hold pending messages. There is a
one-to-one correspondence between the modules of a domain and the mailboxes. In other words,
each module has exactly one mailbox and each mailbox services exactly one module. The main task
of a module is normally blocked trying to read a message from that module's mailbox. When a
message is available, the main task fetches it, processes it, and then (likely) blocks again
trying to read the next message. Since the mailboxes have, in general, space for holding
multiple messages, it is possible that a message can arrive while the module was processing an
earlier message. In that case, the new message is retrieved at once. Modules never process more
than one message at a time. Outstanding messages are queued in the module's mailbox. The mailbox
array thus defines the domain and forms the communication infrastructure for the domain.

In normal operation, all module tasks are blocked, and the CPU is idle. When a hardware
interrupt is generated, an interrupt service routine inside one a driver module for the
associated hardware device activates and sends a message, as appropriate, to another,
higher-level module to handle the interrupt. That module may send other messages, etc., creating
a cascade of activity in the system. Once the interrupt has been fully handled and that activity
dies down, all modules return to their quiescent state of being blocked waiting for a message.
If multiple interrupts happen at about the same time, this presents no problem as the modules
are able to queue messages and execute in parallel as needed. However, the modules do need to be
coded with this scenario in mind. CubedOS is thus a \newterm{reactive} programming system since
it reacts to external input, but does not typically initiate any activity.

CubedOS comes with a collection of \newterm{core modules} that provide services commonly needed
by many applications. The publish/subscribe server mentioned earlier is an example of a core
module. There is also a timeserver module that can be requested to send messages to other
modules at specified times or periodically. "Reacting" to such messages gives modules a way to
perform activities without other specific hardware stimulation (i.e., routine monitoring,
telemetry reporting, and so forth).

CubedOS also comes with a collection of \newterm{library packages} that contain general-purpose
reusable code of various kinds. Unlike a module, a library package does \emph{not} contain a
task\footnote{In a previous version of CubedOS, library packages were called ``passive modules''
to distinguish them from active modules with tasks. That terminology has been dropped.}.
Subprograms inside library packages are called in the usual way and return results in the usual
way, and not via messages\footnote{It is permitted for a library package to send a message as a
side effect.}.

Library packages can also be shared by modules. One consequence of this is that library packages
can have no modifiable global data\footnote{Unless wrapped in a protected object.} since they
might be called by multiple tasks at once. All task interactions are via the message passing
mechanism. Modules should not attempt to communicate via global data inside library packages.

The rest of this document walks you through the steps for creating a simple, single-domain
CubedOS application. You can use these steps as a framework for creating your own applications.
Although we build the application here incrementally for educational purposes, you can review
the final form of the application in the \filename{moonshot} folder of the \filename{samples}
folder in the repository.

\subsection{Moonshot Application}
\label{sec:manual-tutorial-moonshot}

The application we will build demonstrates many, but not all, aspects of CubedOS development. It
is intended to be easy to understand, implement, and modify. It is not intended to be an
exhaustive demonstration of CubedOS and its libraries.

We will create an application, called \newterm{Moonshot}, that roughly simulates a CubeSat
mission orbiting the Moon. It is not our intention for this application to be an accurate
simulation of a real mission. Many details will be left out or simplified. However, imagining
that the application might be controlling a real spacecraft will hopefully make it more fun to
think about!

Although in a real mission, the application would obviously need to control real hardware, in
this tutorial we will simulate the hardware. The modules that might ordinarily serve as device
drivers will instead implement ``mock'' devices with simulated behavior. This allows the
application to run on any development system you might have available.

Moonshot will make use of several CubedOS core modules. Specifically:
\begin{enumerate}
  \item The Log Server will be used to log messages of various kinds for later transmission to
  the Earth. The Log Server is widely used. Almost all CubedOS applications will require it.

  \item The Time Server will be used to trigger periodic actions in other modules. In Moonshot,
  without any real hardware to generate interrupts, the Time Server will be the initiator of all
  activity once system initialization is complete.

  \item The File Server will be used to store image files taken by the mock camera. The Moonshot
  mission will entail taking picture of the lunar surface, so we need a place to store those
  images. In a real mission it is common for there to be some sort of non-volatile storage
  precisely for this purpose.

  \item The Publish/Subscribe Server will be used to allow modules to subscribe to various
  channels. The point of the Publish/Subscribe Server is to decouple senders and receivers of
  messages so that reusable modules can be more readily created. We describe this in more detail
  later.
\end{enumerate}

Moonshot will also have several application-specific modules for managing the (mock) hardware
devices. Specifically:
\begin{enumerate}
  \item The Camera module will simulate a camera that can take pictures of the lunar surface.

  \item The Radio module will simulate a radio that can send messages to and receive messages
  from the Earth.

  \item A Thruster module will simulate maneuvering thrusters that can adjust the orbit and
  orientation of the CubeSat. 

  \item A Controller module that will server as the main control logic for the mission. It will
  coordinate the action of the other modules to fulfill the mission objectives.
\end{enumerate}

To get started, you will need space for your application's files. You will also need to clone
the CubedOS repository from GitHub. At the time of this writing CubedOS is distributed in source
form via its development repository. There is no release distribution yet.

The folder for your application need not be included in the CubedOS repository folder. In fact,
we recommend that you do \emph{not} include it there. Instead, create a folder somewhere else on
your system for your application. You can organize this folder as you please, but for this
example, we will assume all the source files are stored directly in the folder.

It is often useful to use an existing application as a baseline for new work. The Echo
sample in the CubedOS repository is a good place to start. First, create a project file in
your application folder named \texttt{moonshot.gpr} that contains the following:
\begin{verbatim}
project Moonshot is
   for Main use ("main.adb");
   for Object_Dir use "build";
   for Source_Dirs use (".", "../../src/modules", "../../src/library");
   for Languages use ("Ada");

   package Compiler is
      for Default_Switches ("Ada")
        use ("-gnat2022",     -- Use Ada 2022
             "-gnatW8",       -- Enable UTF-8 source files
             "-fstack-check", -- Enable stack checking
             "-gnatwa",       -- Enable all warnings
             "-gnata",        -- Enable assertions
             "-g");           -- Enable debugging
   end Compiler;
end Moonshot;
\end{verbatim}

The source paths should include the current folder (``.''), or wherever you plan to store the
source files for your application, together with the two source folders from the CubedOS
repository. Edit the paths above as appropriate for your system. Create a folder named
\texttt{build} in your application folder (so the relative path in the project file works) to
hold the build artifacts of the application. The default switches are fine for this getting
started application. They enable all warnings, runtime assertion checking, and debugging
support. In a real application you may want different switches or multiple build scenarios.

Notice that CubedOS applications are intended to use Ada 2022 as their base language. Although
this is not strictly necessary, Ada 2022 does provide some niceties that are helpful. Finally,
the |-gnatW8| switch enables support for Unicode UTF-8 encoded source files, allowing certain
Unicode characters, such as Greek letters, in identifier names. That is also not strictly
necessary, but it can be helpful for clarity in some cases.

Every Ada program requires a library-level, parameter-less procedure to serve as the program's
entry point. CubedOS applications don't actually use this procedure, but it is required for the
Ada language. Use the name \texttt{main.adb} as indicated above. We will describe what must be
in this file later.

Filling in the rest of the application requires the following steps:
\begin{enumerate}
\item Instantiating the CubedOS generic message manager to create a domain (mailbox array) for
  your application.
\item Writing a name resolver package to assign addresses to the modules you will use, including
  the core modules.
\item Writing a publish/subscribe channel definition package to define the channels you will use
  in your application.
\item Writing MXDR interface definition descriptions for the messages that will be sent between
  modules, and using the Mercury tool to generate the Ada code for those message encoding and
  decoding subprograms.
\item Creating skeletons for the modules you will write by copying the templates from the
  CubedOS repository and editing them.
\item Writing the Ada-required main procedure to ensure the modules you need will be compiled
  into your executable.
\item Verifying that your completed skeleton application builds and runs.
\item Implementing the full logic of your application by filling in the skeletons you created
  earlier.
\end{enumerate}
